\section{The common history and its centered service}
\label{sec:selected-centered-service}
We now apply the local constructions to a single family. The input is
the simultaneous purchase of
Sections~\ref{sec:common-funded-history} and
\ref{sec:common-observable-purchases}, including the stationary and
free convergence estimates on their common horizon. We retain the
prescribed positive blocking penalty throughout.

Choose nondecreasing profile orders and spatial powers, both at least the
stage index, and nondecreasing support exponents at least $100$. The
particular choices $6+j$, $2+j$ and $100+j$ suffice for the elementary
construction. The selected public run uses the first stage $4$, the
orders $6+j$, the powers $2+j$, and the constant support exponent
$b_j=100$ at every stage; the identity $b_j=100$ is retained with that
run and is required by later consumers. Throughout this section the
public run $R$ is retained, including its selected shooting parameters,
its reference schedule $t_R$, its transition records and its transfer
family. The following paragraphs unpack the construction and the
estimates of these retained witnesses; where they speak of choosing an
object, the object meant is the one stored in $R$, and no second choice
is made. Choose the fixed parameter radius and the small response reserve
as in Section~\ref{sec:common-observable-purchases}. Make the block-length
choices before their numerical costs. Apply
Section~\ref{sec:common-scale-selection} to all the resulting requests,
with a nondecreasing occupancy schedule. This gives one sequence of
levels, couplings and periods, with $J(m)\to\infty$.

\subsection{Choosing the history before its observable comparisons}
For every cutoff $N$ and period $p$, the terminal family is the exact
Wilson family of Section~\ref{sec:wilson-startup-construction}.
The interval return and scalar reset are applied with the strengthened
local conditions of Section~\ref{sec:common-observable-purchases}.
Thus every member of the resulting continuous interval family carries
those conditions. To prove that its bottom couplings cover the required
interval, forget only these additional conditions and apply the barrier
and covering argument of Section~\ref{sec:surviving-interval-shooting}.
This leaves the interval, its endpoints and every replayed history
unchanged. Choose a parameter in that very interval whose bottom
coupling is the prescribed common value, and retain its original
strengthened certificate.
Retain that member, its actual representations $R_{N,p,m}$ and its
adjacent transitions. All of them use the already purchased paths,
kernels and clipping data. Choosing the members for the countably
many pairs $(N,p)$ introduces no new scale request. On the run $R$ these
are stored data: for every $(N,p)$ the run stores the interval family, a
parameter in its interval, and the identity of its history at $(N,p)$
with the replay of that family at that parameter. Its nesting receipt
records that every adjacent pair of this history is a protected and
accurate transition of the purchased step and that the terminal
representative carries the reserved addresses.

The original certificate at the chosen parameter supplies the local
parameter estimates on these histories. Thus the comparisons and fixed-level limits in
Sections~\ref{sec:coupled-background-convergence},
\ref{sec:selected-history-limits} and
\ref{sec:fixed-rung-convergence} apply to this family. The uniform
local derivative bank is that of
Section~\ref{sec:common-observable-purchases}. In particular, the
fixed-level convergence is available for the actual selected
activity arrays, not just for their stationary paths.

This application carries one input that the certificate of the
selected member does not by itself provide.
Section~\ref{sec:coupled-background-convergence} compares the selected
history with auxiliary truncated histories, which have a prescribed
coupling script and zero terminal activity, and
Section~\ref{sec:selected-history-limits} draws its limits under the
assumption that the transition hypotheses of
Section~\ref{sec:parameter-comparison-suppliers} hold for the actual
input of each such auxiliary history at every step, on the same path
and dictionaries. The certificate at the chosen parameter concerns
the selected history only. The fixed-level limits are therefore used
here subject to that same assumption on the auxiliary histories; it
is not proved in this section.

Put $t_d=t_R(d)$, the reference level stored in the run for the grade
$d$. The public construction supplies
\[
                 1\le t_d,\qquad d+3\le J(t_d).
\]
Such levels exist by divergence of $J$; no monotonicity of
$d\mapsto t_d$ is recorded or used. At each adjacent step
choose the wing containing the origin principal modes and the keeper
keys for all relevant support/inventory shapes. These keys form a finite
set: the torus is finite, so its support powerset and the pairs
$I\subseteq X$ with $I\ne\varnothing$ are finite. The wing uses
these shapes, not the countably many reserved serial addresses. Hence
it covers every forced coefficient while remaining finite. Take the
finite union required for the admitted grades.
Section~\ref{sec:heading-history-handoff} supplies one heading and
one local transfer map for that entire wing. Consequently the same
map works for all those grades. Enlarging the wing preserves its
finite interpolation proof and the exact endpoint representation.

Apply the finite backward construction of
Section~\ref{ex-centered-freight-producer}, starting at the terminal
marked section constructed in Section~\ref{sec:terminal-marked-construction}. This gives the finite transfer data at every grade,
cutoff and period. The three exact identities in
\eqref{eq:history-three-identities} identify the coefficient,
exponential entry and forced profile throughout the construction.
In particular, its terminal density is
$a_N^d R_{N,p,N}(U)\,\mathrm{initial}_{N,p}(v)(U)$;
the physical observable $\mathrm{initial}_{N,p}$ already contains
the factor $a_N^{-d}$. The local maps used
in this construction are the adjacent maps already chosen above;
the downward induction does not select another background history.
The unit, gauge invariance and signed-motion identities are preserved
by the exact local section and by the equivariant finite assembly.
Thus these transfer data form one family on the retained history. On the
run $R$ this family is the stored transfer family $S_R$, and the adjacent
maps are its stored transition records; for $t_d>N$ the stored transfer
datum is a default map and carries no assertion, as explained in the
proof of Proposition~\ref{prop:ex-original-history-attachment}.

\subsection{The service at every reference level}
Fix a grade $d$, a period $p$, a cutoff $N$, and
$t_d\le k\le N$; in particular $k\ge1$. We verify the inputs of
Theorem~\ref{thm:ex-selected-run-service} for this family.

At every step above $k$, the ordinary and raised menu estimates of
Section~\ref{sec:finite-menu-history-handoff} apply to the selected
transition. Its input and output are exactly $R_{N,p,m+1}$ and
$R_{N,p,m}$. Its three tariff inequalities give the local duty
bound and the same single-marker derivative. The terminal clause
is supplied by Proposition~\ref{prop:terminal-marked-voice}:
the marked density and forced-profile identity are
constructed on the exact startup representation, from the reserved
addresses recorded in the nesting receipt of the run, with their separate
entrance fee paid before scale selection
(Section~\ref{sec:common-observable-purchases}).
This proves retention all the way to the terminal level.

The room and bank constants are those chosen for the local
coefficient maps. The scale inequalities in
Section~\ref{ex-centered-geyser-scale-estimates} give their common
marker prices
\[
 \zeta_{d,m}=\widehat\beta_m^{-(\deg F_d+4)},\qquad
 \phi_m=\widehat\beta_{m+1}^{-1},\qquad
 \delta_m=\widehat\beta_{m+1}^{-2},\qquad
 \eta_m=\tfrac14\widehat\beta_m^{-1/2}.
\]
Here $F_d$ is the finite coefficient frame. Their strict positivity,
local coefficient bounds, adjacent-price growth and support inequalities
are proved there. Since $d\le J(k)\le J(m)$, the finite-grade
inequalities purchased at stage $J(m)$ apply at every $m\ge k$.
The occupancy bound at level $m$ is $N_{J(m+1)}$, which is
nondecreasing because both maps are nondecreasing.

The same scale estimates give
\[
 B_m\ge8,\qquad
 \delta_m\phi_m e^{B_m/128}\ge1,\qquad
 B_{m+1}^2\le L_mB_m^2.
\]
The stain/contact bound pays the simultaneous forced support
requirement. These inequalities continue to hold after restriction
to any reference level $k$; no constant is selected from $N$.
For the local-section norm use its fixed camera constant $C_d$.
The fitted stencil has cover at most one and its values are bounded
by $C_d$, uniformly in the cutoff and period. Its paid coefficient
price is precisely the bank price in the scale estimates. The
cover bound just stated is the footprint-cover clause of the height
record, item~(iii) of Theorem~\ref{thm:ex-selected-run-service}; the
support-size inequalities of its item~(ii) are among the scale
estimates displayed above. Section~\ref{ex-centered-site-geometry}
is not used for either: it supplies the occupancy, migration and
separation of the label sites, which enter only in the moment
bound below.

The backward transfer data supply the remaining finite source
record. Its coefficient and activity bounds are the two separate
ledgers in Section~\ref{ex-centered-freight-producer}; its adjacent
identities use the selected maps, and its terminal forced-profile
identity is the third identity in
\eqref{eq:history-three-identities}. That identity is essential:
the scalar integral identity alone would not give the terminal
marked polynomial. All these conclusions hold for every $(d,k,N,p)$
in the stated range, with the already fixed reference levels.
They are exactly the service hypotheses, now supplied by the
local purchases and their finite backward construction.

We may therefore apply Theorem~\ref{thm:ex-selected-run-service}
and the terminal birth calculation. At a reference level $k$, a
finite label list with paid weights and separated origins has
normalized moment bound
\[
              \exp\{|\Lambda_k|\mathcal T_k+6|\mathcal V|\}.
\]
Let $m\le k$ now denote the root level opened at $k$, as in
Section~\ref{ex-centered-site-geometry}, rather than a running level.
Take one label at each site of $\Lambda_m$, so that
$M=|\mathcal V|=|\Lambda_m|$, placed by the aligned embeddings of
Section~\ref{ex-centered-site-geometry}. With $n_j$ the side of the
periodic lattice $\Lambda_j$ and $q_k=a_m/a_k$, that section records
$n_k=q_kn_m$; since $|\Lambda_j|=n_j^4$, this gives the site identity
$|\Lambda_k|=M(a_m/a_k)^4$.
The normalized $M$th root is consequently bounded by
\[
              \exp\{(a_m/a_k)^4\mathcal T_k+6\}/w.
\]
This is the centered estimate for the same family, simultaneously
at all suitable periods. The fixed-level convergence just retained
and the removal of higher marker terms in
Section~\ref{sec:fullmesh-correlation-cauchy} then give convergence
along the full mesh sequence at every fixed suitable period.
The finite-difference bounds of
Section~\ref{sec:ex-fullmesh-fixedbounds} apply to the same
transfer data, so its limiting correlations have the smoothness
proved in Section~\ref{ex-smooth-limits}.

The further calibration of an arbitrary marked section and the
rank and ultraviolet comparisons must still be made on this family.
They will use the same reference levels and transfers; the centered
construction above is already fixed before those comparisons or a
volume subsequence are chosen.

\paragraph{Relation to the paid selected-history record.}
Section~\ref{sec:actual-marked-portfolio-calibration} consumes a record
(R1)--(R7) on the same family and the same transition records. Its
parameters are those of this section: the stage function $W$ is the
purchased worth, $u_i=\eta_i$, $\phi_i$, $\delta_i$ and $\zeta_{d,i}$
are the marker prices displayed above, $r_i=N_{J(i+1)}$, the doors are
the stored transition records, and $\psi_i$ is the common waybill
\eqref{eq:history-common-waybill} at level $i+1$. With these
parameters, the present section assembles the following items, each from
the inputs named in its place above. The room and bank inequalities (R1)
are those of the purchased budget. The scale bounds (R2) are the
positivity of the marker prices and the three displayed inequalities.
The monotonicity in (R3) is that of the occupancy bound. The retention
clauses and the terminal clause of the second paragraph of this
subsection are (R4). Of (R5), this section supplies $\phi_i\le1$,
$\psi_i\ge0$, the second, third and fourth displayed inequalities, the
bound $X=C_d$, the cover number one, and the terminal inequality for
$u_N$. The backward transfer data, with their door-by-door comparison at
price $\psi_i$ and their terminal forced-profile identity, are (R6).
Three items are not part of the centered service and are not derived in
this section: $r_i\ge1$ in (R3); the positivity $\zeta_{d,i}>0$ and the
first inequality of (R5), which is
\eqref{eq:ex-marked-height-jump} for the prices above; and the paid
one-step input (R7). They belong to the stronger form of the funded
continuation described in the proof of
Proposition~\ref{prop:ex-original-history-attachment}: $r_i\ge1$ is the
prescribed lower bound of the occupancy schedule, positivity is that of
$\widehat\beta_i$, the first inequality of (R5) is the tail clause of
the delay purchased in Section~\ref{sec:common-observable-purchases},
and (R7) is the local transition estimate obtained in that proof from
the nesting receipt, the transition records and the floor-length
condition of Section~\ref{sec:common-observable-purchases} at grade
zero. For none of the items has this
section a self-contained proof: each rests on the cited producers.
